\documentclass[11pt]{article}
\usepackage[a4paper,margin=1in]{geometry}
\usepackage{fontspec}
\usepackage[boldfont=false]{xeCJK}
\setCJKmainfont{FandolSong-Regular.otf}
\usepackage{amsmath}
\usepackage{amssymb}
\usepackage{array}
\usepackage{booktabs}
\usepackage{graphicx}
\usepackage{adjustbox}
\usepackage{enumitem}
\setlist[itemize]{topsep=2pt,partopsep=0pt,parsep=0pt,itemsep=2pt,leftmargin=1.4em}
\setlist[enumerate]{topsep=2pt,partopsep=0pt,parsep=0pt,itemsep=2pt,leftmargin=1.6em}
\usepackage{parskip}
\usepackage{fvextra}
\fvset{breaklines=true,breakanywhere=true,fontsize=\small}
\setlength{\parindent}{0pt}

\begin{document}

\begin{center}
  {\LARGE\bfseries Stoichiometry}\\[0.35em]
  {\large Igcse Chemistry}
\end{center}
\vspace{0.6em}

\textbf{Stoichiometry} 化学计量 is the study of the amounts of substances in a reaction — how much reacts and how much is made. This topic is mostly about counting atoms and doing calculations.

\section*{Chemical formulae}

A \textbf{formula} 化学式 shows which \textbf{atoms} 原子 are in a substance, and how many of each. There are two kinds of formula you must know.

\begin{itemize}
\item The \textbf{molecular formula} 分子式 is the actual number of each atom in one \textbf{molecule} 分子. For glucose it is $\text{C}_6\text{H}_{12}\text{O}_6$.

\item The \textbf{empirical formula} 实验式 is the simplest whole-number \textbf{ratio} 比例 of the atoms or \textbf{ions} 离子 in a \textbf{compound} 化合物. For glucose it is $\text{CH}_2\text{O}$.

\end{itemize}

\subsection*{Working out a formula}

If you are given a model or diagram, just count the atoms of each element and write them as a formula.

For an \textbf{ionic compound} 离子化合物 you can work out the formula from the charges on the ions. The total positive charge must balance the total negative charge, because the compound has no overall charge. Some common ions:

\begin{center}
\adjustbox{max width=\linewidth}{%
\begin{tabular}{ll}
\toprule
\textbf{Positive ions} & \textbf{Negative ions} \\
\midrule
$\text{Na}^{+}$, $\text{K}^{+}$, $\text{H}^{+}$ & $\text{Cl}^{-}$, $\text{OH}^{-}$, $\text{NO}_3^{-}$ \\
$\text{Mg}^{2+}$, $\text{Ca}^{2+}$, $\text{Cu}^{2+}$ & $\text{O}^{2-}$, $\text{SO}_4^{2-}$, $\text{CO}_3^{2-}$ \\
$\text{Al}^{3+}$ & $\text{N}^{3-}$ \\
\bottomrule
\end{tabular}}
\end{center}

For example, $\text{Na}^{+}$ and $\text{O}^{2-}$: you need two $\text{Na}^{+}$ to balance one $\text{O}^{2-}$, so the formula is $\text{Na}_2\text{O}$.

\section*{Writing equations}

An equation shows how \textbf{reactants} 反应物 (the starting substances) change into \textbf{products} 生成物 (the substances made).

\begin{itemize}
\item A \textbf{word equation} 文字方程式 uses names: \emph{magnesium + oxygen → magnesium oxide}.

\item A \textbf{symbol equation} 化学方程式 uses formulae and must be \textbf{balanced} 配平 — the same number of each atom on both sides.

\end{itemize}

\[
2\text{Mg} + \text{O}_2 \rightarrow 2\text{MgO}
\]

You add \textbf{state symbols} 状态符号 in brackets to show the physical state: $(s)$ solid, $(l)$ liquid, $(g)$ gas, and $(aq)$ for \textbf{aqueous} 水溶液 (dissolved in water).

\[
\text{Zn}(s) + 2\text{HCl}(aq) \rightarrow \text{ZnCl}_2(aq) + \text{H}_2(g)
\]

\subsection*{Ionic equations}

An \textbf{ionic equation} 离子方程式 shows only the ions that actually change. Ions that are the same on both sides are \textbf{spectator ions} 旁观离子 and are left out. For example, when an acid reacts with an alkali:

\[
\text{H}^{+}(aq) + \text{OH}^{-}(aq) \rightarrow \text{H}_2\text{O}(l)
\]

\section*{Relative masses}

The \textbf{relative atomic mass} 相对原子质量 ($A_r$) of an \textbf{element} 元素 is the average \textbf{mass} 质量 of its atoms compared to $\tfrac{1}{12}$ of the mass of one $^{12}\text{C}$ atom.

The \textbf{relative molecular mass} 相对分子质量 ($M_r$) is the sum of the relative atomic masses of all the atoms in the molecule. For ionic compounds we use the \textbf{relative formula mass} 相对式量, found the same way from the formula.

\[
M_r(\text{H}_2\text{O}) = (2 \times 1) + 16 = 18
\]

\subsection*{Reacting masses by simple proportion}

You can sometimes find a reacting mass without the mole. If 24 g of magnesium makes 40 g of magnesium oxide, then 12 g of magnesium (half as much) makes 20 g of magnesium oxide.

\section*{The mole}

The \textbf{mole} 摩尔 (symbol mol) is the unit for the \textbf{amount of substance} 物质的量. One mole of any substance contains $6.02 \times 10^{23}$ \textbf{particles} 粒子 (atoms, ions or molecules). This number is the \textbf{Avogadro constant} 阿伏伽德罗常数.

The \textbf{molar mass} 摩尔质量 is the mass of one mole, in grams per mole (g/mol). Its number is the same as the $A_r$ or $M_r$. The key relationship is:

\[
\text{amount (mol)} = \frac{\text{mass (g)}}{\text{molar mass (g/mol)}}
\]

\textbf{Worked example.} How many moles are in 36 g of water? Molar mass of water $= 18$ g/mol.

\[
n = \frac{36}{18} = 2 \text{ mol}
\]

To find the number of particles, multiply the moles by the Avogadro constant.

\subsection*{Volumes of gases}

At room temperature and pressure (r.t.p.), one mole of any gas takes up the same \textbf{volume} 体积. This \textbf{molar gas volume} 摩尔气体体积 is $24 \text{ dm}^3$.

\[
\text{volume of gas (dm}^3) = \text{amount (mol)} \times 24
\]

\subsection*{Concentration of solutions}

The \textbf{concentration} 浓度 of a \textbf{solution} 溶液 can be given in $\text{g/dm}^3$ or in $\text{mol/dm}^3$.

\[
\text{concentration (mol/dm}^3) = \frac{\text{amount of solute (mol)}}{\text{volume (dm}^3)}
\]

Remember to convert volume: $1 \text{ dm}^3 = 1000 \text{ cm}^3$, so divide a volume in $\text{cm}^3$ by 1000.

\section*{Doing reaction calculations}

Most calculations follow the same steps: change the known amount into moles, use the balanced equation to find the moles of what you want, then change back into mass, volume or concentration.

\subsection*{Limiting reactant}

When two reactants are mixed, often one runs out first. This is the \textbf{limiting reactant} 限量反应物. It decides how much product can form; any other reactant is in excess (left over).

\subsection*{Titration calculation}

In a \textbf{titration} 滴定 you measure the volume of one solution that reacts with another. From the volume and concentration you find the moles of one \textbf{solute} 溶质, then use the equation ratio to find the moles, concentration or volume of the other.

\subsection*{Empirical and molecular formulae from data}

To find an empirical formula from masses or percentages:

\begin{enumerate}
\item Divide each element's mass (or \%) by its $A_r$.

\item Divide all the answers by the smallest one to get the simplest ratio.

\end{enumerate}

To find the molecular formula, compare the empirical formula mass with the real $M_r$ and multiply up.

\subsection*{Percentages}

\begin{itemize}
\item \textbf{Percentage yield} 产率 compares how much product you actually got with the most you could get: $\dfrac{\text{actual}}{\text{theoretical}} \times 100$.

\item \textbf{Percentage composition by mass} 质量分数 of an element $= \dfrac{\text{mass of that element in the formula}}{M_r} \times 100$.

\item \textbf{Percentage purity} 纯度 $= \dfrac{\text{mass of pure substance}}{\text{mass of impure sample}} \times 100$.

\end{itemize}


\end{document}
